As for the probability, supposing that the neutron star space orientation is
random, is
\[
p = \frac{2A_p(d)}{4\pi d^2} \tag*{(2 pt)}
\]
where $A_p(d)$ is the area of the path of one light cone.

% FIGURE NEEDED: hemisphere of radius d with the double emission cone traced on
% it -- the beam path is an annular band between two spherical caps of heights
% D and H, half-angles 28 deg and 32 deg (28+4) marked (2012_th_sol.pdf p. S3).

As the figure above shows, the area is the difference of the spherical caps of
height $H$ and $D$:
\[
A = 2\pi d (H - D)
  = 2\pi d^2 \left[\cos\left(\theta - \frac{\alpha}{2}\right)
    - \cos\left(\theta + \frac{\alpha}{2}\right)\right] \tag*{(3 pt)}
\]
\[
p = \cos\left(\theta - \frac{\alpha}{2}\right)
  - \cos\left(\theta + \frac{\alpha}{2}\right)
  = 3.5 \cdot 10^{-2} \quad \text{or} \quad 3.5\% \tag*{(2 pt)}
\]
The light waves are spread along the double cone and not along a sphere, thus:
\[
F = \frac{L}{2A(d)}
  = \frac{10\,000\,L_\odot}{4\pi d^2\left(1 - \cos\dfrac{\alpha}{2}\right)}
\]
Where $A(d)$ is the area of one light cone. Using the Sun's absolute magnitude,
we can find the absolute magnitude of the pulsar ($m_p$):
\[
m_p - M_\odot
  = -2.5\log\left(\frac{F}{\dfrac{L_\odot}{4\pi(10~\mathrm{pc})^2}}\right)
  = -2.5\log\left(10\,000 \cdot \frac{1}{1 - \cos\dfrac{\alpha}{2}}
    \cdot \left(\frac{10}{1000}\right)^2\right) \tag*{(3 pt)}
\]
\[
m_p = -3.24
\]
